% vim: set tw=78 tabstop=4 shiftwidth=4 aw ai:
\documentclass{beamer}

\usepackage[utf8x]{inputenc}		% diacritice
\usepackage[english]{babel}
\usepackage{color}			% highlight
\usepackage{alltt}			% highlight
\usepackage{multicol}
\usepackage{underscore}

% highlight; comment this out in case you don't input code source files
% \usepackage{code/highlight}		% highlight

\usepackage{hyperref}			% folosiți \url{http://...}
					% sau \href{http://...}{Nume Link}
\usepackage{verbatim}
\usepackage{wasysym}
\usepackage{subfigure}

% Show contents at every section beginning. Ripped off from manual.
\AtBeginSection[] % Do nothing for \section*
{
	\begin{frame}<beamer>
		\frametitle{Outline}
	\tableofcontents[currentsection]
		\end{frame}
}

\mode<presentation>
{ \usetheme{Pittsburgh} \usecolortheme{seahorse} }

\setbeamertemplate{footline}[institute]{}

% Încărcăm simbolurilor Unicode românești în titlu și primele pagini
\PreloadUnicodePage{200}

\author{
		{\hspace{3mm} Adriana Costin \hspace{12mm}  Adrian Stanciu}
		\newline
		{Mihai Maruseac \hspace {10mm} Irina Grădinaru \hspace{10mm}
		Adrian Șendroiu}}
\title{Advances in Information Retrieval}

\institute[ACS UPB]{Faculty of Automatic Control and Computer Science \\
					Politehnica University of Bucharest}
\date{January 11, 2013}

\begin{document}

% Slide-urile cu mai multe părți sunt marcate cu textul (cont.)
\setbeamertemplate{frametitle continuation}[from second]

% Arătăm numărul frame-ului
% \setbeamertemplate{footline}[frame number]

\frame{\titlepage}

\frame{\tableofcontents}

% NB: Secțiunile nu sunt marcate vizual, ci doar apar în cuprins
\section{Dummy}

\begin{frame}{Cattitude}
	\begin{columns}
		\begin{column}[l]{0.2\textwidth}
		\begin{itemize}
			\item miau1
			\item miau2
			\item miau3
		\end{itemize}

		\end{column}
		\begin{column}[l]{0.8\textwidth}
			\begin{figure}
				\includegraphics[scale=1.5]{img/cat.jpg}
			\end{figure}
		\end{column}
	\end{columns}
\end{frame}


\section{Trends and Ideas}

\begin{frame}{Trends and Ideas}

Ce-a zis profu în mail:
\begin{quotation}
	 Puneti la final o sectiune cu:
	\begin{itemize}
		\item Trenduri (ce vi se pare voua ca e la moda pe directia aia acum).
		\item Idei (ce credeti ca s-ar putea face pentru a avansa directia respectiva).
	\end{itemize}
	
\end{quotation}

\end{frame}



\end{document}
